% GENERATED FILE. Edit canonical lessons, then run npm run generate:books.
% canonical-source-hash: fnv1a64:4bc172c08315cbe1
% canonical-lessons: ZH-C71-haohe, ZH-C71-tian, ZH-C71-la, ZH-C71-xian, ZH-C71-diancai

\chapter{Tastes}
\label{ch:zh-tastes}
\hlchaptermodality{71}

\begin{chapteropening}
\textbf{By the end of this chapter:} \emph{I can say a drink is nice, food is sweet, hot or salty, and order a meal.}

Complete the last lesson of chapter 71: I can say a drink is nice, food is sweet, hot or salty, and order a meal.
\end{chapteropening}

\section[\texorpdfstring{\zh{好喝} (hǎohē)}{好喝 (hǎohē)}]{\zh{好喝} (hǎohē) — nice to drink}

\label{lesson:ZH-C71-haohe}



\begin{quote}
Before the new one: say the Mandarin for dumpling, then the Mandarin for soup.
\end{quote}

\subsection*{You'll want to know: \zh{好喝}}
\textbf{\zh{好喝}} — \emph{hǎohē} — \textquotedblleft{}nice to drink\textquotedblright{}.

\begin{quote}
\emph{zhèr de chá hěn hǎohē} — The tea here is lovely.
\end{quote}

\begin{cousinweb}
\textbf{\zh{好喝}} (\emph{hǎohē}) is built from \zh{好} 'good' and \zh{喝} 'to drink', both inherited from Middle Chinese, \emph{xawX} and \emph{xat} (Guangyun, in Baxter's transcription); the dictionary glosses the second as 'to shout at', and the drinking sense is not given there.
\end{cousinweb}

\begin{grammarlens}[title={What you've built}]
One more word for everyday answers.
\end{grammarlens}

\subsection*{Your turn}
\begin{itemize}
\raggedright
  \item \emph{Say it:} \emph{hǎohē}
  \item \emph{Say it:} \emph{hǎohē}, once more
  \item \emph{Say it:} \emph{zhèr de chá hěn hǎohē} — The tea here is lovely.
  \item \emph{Recall:} say the Mandarin for dumpling, then the Mandarin for soup, then say \emph{hǎohē} again
\end{itemize}

\begin{tcolorbox}[breakable,colback=teal!4,colframe=teal!35!black,title={Before you move on}]
What does \zh{好喝} mean? (\textquotedblleft{}Nice to drink\textquotedblright{}.) Say it once more.
\end{tcolorbox}
\section[tián]{tián — sweet}

\label{lesson:ZH-C71-tian}



\begin{quote}
Before the new one: say the Mandarin for soup, then the Mandarin for nice to drink.
\end{quote}

\subsection*{You'll want to know: tián}
\emph{tián} — \textquotedblleft{}sweet\textquotedblright{}.

\begin{quote}
\emph{nǐ de kāfēi hěn tián} — Your coffee is very sweet.
\end{quote}

Its characters come later; for now it is a word for the ear.

\begin{cousinweb}
\emph{tián} is inherited from Middle Chinese \emph{dem} (Guangyun, in Baxter's transcription), glossed 'sweet'.
\end{cousinweb}

\begin{grammarlens}[title={What you've built}]
One more word for everyday answers.
\end{grammarlens}

\subsection*{Your turn}
\begin{itemize}
\raggedright
  \item \emph{Say it:} \emph{tián}
  \item \emph{Say it:} \emph{tián}, once more
  \item \emph{Say it:} \emph{nǐ de kāfēi hěn tián} — Your coffee is very sweet.
  \item \emph{Recall:} say the Mandarin for soup, then the Mandarin for nice to drink, then say \emph{tián} again
\end{itemize}

\begin{tcolorbox}[breakable,colback=teal!4,colframe=teal!35!black,title={Before you move on}]
What does tián mean? (\textquotedblleft{}Sweet\textquotedblright{}.) Say it once more.
\end{tcolorbox}
\section[là]{là — spicy; hot}

\label{lesson:ZH-C71-la}



\begin{quote}
Before the new one: say the Mandarin for nice to drink, then the Mandarin for sweet.
\end{quote}

\subsection*{You'll want to know: là}
\emph{là} — \textquotedblleft{}spicy; hot\textquotedblright{}.

\begin{quote}
\emph{zhèr de cài hěn là} — The food here is really spicy.
\end{quote}

Its characters come later; for now it is a word for the ear.

\begin{cousinweb}
\emph{là} is inherited from Middle Chinese \emph{lat} (Guangyun, in Baxter's transcription), glossed 'pungent'; the dictionary writes it with a variant form of the character.
\end{cousinweb}

\begin{grammarlens}[title={What you've built}]
One more word for everyday answers.
\end{grammarlens}

\subsection*{Your turn}
\begin{itemize}
\raggedright
  \item \emph{Say it:} \emph{là}
  \item \emph{Say it:} \emph{là}, once more
  \item \emph{Say it:} \emph{zhèr de cài hěn là} — The food here is really spicy.
  \item \emph{Recall:} say the Mandarin for nice to drink, then the Mandarin for sweet, then say \emph{là} again
\end{itemize}

\begin{tcolorbox}[breakable,colback=teal!4,colframe=teal!35!black,title={Before you move on}]
What does là mean? (\textquotedblleft{}Spicy; hot\textquotedblright{}.) Say it once more.
\end{tcolorbox}
\section[xián]{xián — salty}

\label{lesson:ZH-C71-xian}



\begin{quote}
Before the new one: say the Mandarin for sweet, then the Mandarin for spicy.
\end{quote}

\subsection*{You'll want to know: xián}
\emph{xián} — \textquotedblleft{}salty\textquotedblright{}.

\begin{quote}
\emph{jīntiān de tāng hěn xián} — Today's soup is salty.
\end{quote}

Its characters come later; for now it is a word for the ear.

\begin{cousinweb}
\emph{xián} is inherited from Middle Chinese \emph{heam} (Guangyun, in Baxter's transcription), glossed 'not bland'; the dictionary writes it with the traditional form of the character.
\end{cousinweb}

\begin{grammarlens}[title={What you've built}]
One more word for everyday answers.
\end{grammarlens}

\subsection*{Your turn}
\begin{itemize}
\raggedright
  \item \emph{Say it:} \emph{xián}
  \item \emph{Say it:} \emph{xián}, once more
  \item \emph{Say it:} \emph{jīntiān de tāng hěn xián} — Today's soup is salty.
  \item \emph{Recall:} say the Mandarin for sweet, then the Mandarin for spicy, then say \emph{xián} again
\end{itemize}

\begin{tcolorbox}[breakable,colback=teal!4,colframe=teal!35!black,title={Before you move on}]
What does xián mean? (\textquotedblleft{}Salty\textquotedblright{}.) Say it once more.
\end{tcolorbox}
\section[diǎncài]{diǎncài — to order food}

\label{lesson:ZH-C71-diancai}



\begin{quote}
Before the new one: say the Mandarin for spicy, then the Mandarin for salty.
\end{quote}

\subsection*{You'll want to know: diǎncài}
\emph{diǎncài} — \textquotedblleft{}to order food\textquotedblright{}.

\begin{quote}
\emph{wǒ xiǎng diǎncài} — I'd like to order.
\end{quote}

Its characters come later; for now it is a word for the ear.

\begin{cousinweb}
\emph{diǎncài} is built from \zh{点} 'a dot, to mark' and cài 'dish', both inherited from Middle Chinese, \emph{temX} and \emph{tshojH} (Guangyun, in Baxter's transcription); the dictionary glosses the second as any plant fit to eat. Marking the dishes you want is the likely picture, but the step is not settled.
\end{cousinweb}

\begin{grammarlens}[title={What you've built}]
One more word for everyday answers.
\end{grammarlens}

\subsection*{Your turn}
\begin{itemize}
\raggedright
  \item \emph{Say it:} \emph{diǎncài}
  \item \emph{Say it:} \emph{diǎncài}, once more
  \item \emph{Say it:} \emph{wǒ xiǎng diǎncài} — I'd like to order.
  \item \emph{Recall:} say the Mandarin for spicy, then the Mandarin for salty, then say \emph{diǎncài} again
\end{itemize}

\begin{tcolorbox}[breakable,colback=teal!4,colframe=teal!35!black,title={Before you move on}]
What does diǎncài mean? (\textquotedblleft{}To order food\textquotedblright{}.) Say it once more.
\end{tcolorbox}
